\documentclass[11pt]{article}
\usepackage[margin=1in]{geometry}
\usepackage{amsmath,amssymb,amsthm}
\usepackage{booktabs}
\usepackage[numbers,sort&compress]{natbib}
\usepackage{hyperref}
\usepackage{xcolor}

\newtheorem{theorem}{Theorem}
\newtheorem{proposition}{Proposition}
\newtheorem{definition}{Definition}
\newcommand{\todo}[1]{\textcolor{red}{[TODO: #1]}}

\title{Bounding the Probability That an LLM Agent Exceeds Its Budget:\\
Run-Level Risk Control with Learn-then-Test}
\author{Robat Das\\
Independent Researcher\\
ORCID: \href{https://orcid.org/0009-0006-7792-8351}{0009-0006-7792-8351}}
\date{DRAFT --- \today}

\begin{document}
\maketitle

\begin{abstract}
LLM agents are billed per token, and the cost of a run is not known until it
ends: the number of steps, the model chosen at each step, and the length of
every response are random. Operators therefore face a choice between reserving
for the expected cost, which overruns the budget an uncontrolled fraction of the
time, and reserving for the worst case, which is safe but so conservative that a
capable model is almost never affordable. We ask for something in between: a
controller whose probability of exceeding a run budget $B$ is at most a
user-chosen $\varepsilon$. We first show why the natural approach fails.
Bounding each step's output length with a conformal or binomial quantile and
splitting $\varepsilon$ across steps is valid but, at realistic calibration
sizes, collapses to worst-case reservation, because the union bound charges
every step for a tail event that almost never causes a run-level overrun. We
then calibrate the controller at the level of whole runs. A one-parameter
family of reservation policies is tested in a fixed sequence with exact
binomial $p$-values (Learn-then-Test), which yields
$\Pr\bigl(\Pr(\text{spend} > B) \le \varepsilon\bigr) \ge 1-\delta$ from
i.i.d.\ runs alone, with no assumption on per-step output distributions. In
simulation the calibrated controller keeps overruns below $\varepsilon$ (in the
no-truncation mode) while
assigning 31--45\% of steps to the capable model, where per-step bounds and
worst-case reservation assign none; across 500 independent calibrations the
guarantee held in 499 (well-specified simulator) and 495 (misspecified
simulator) cases at $\delta = 0.05$. The method is implemented in the
open-source library Baar-Core. On 2{,}100 real gpt-4o-mini and gpt-4o agent
runs on HotpotQA, replayed exactly under each policy, worst-case reservation
finished no run; the calibrated policy matched expected-cost reservation's
completion at looser budgets, traded some completion for fewer overruns at
the tightest one, and its guarantee held in at least 499 of 500
calibrations in every setting.
\end{abstract}

% ---------------------------------------------------------------------------
\section{Introduction}

An LLM agent alternates between model calls and tool calls until it produces
an answer. Each model call costs
$p^{\text{in}}_m n_t + p^{\text{out}}_m y_t$, where $m$ is the model, $n_t$ the
prompt length (which grows as the conversation grows) and $y_t$ the output
length. Only $n_t$ is known before the call. Measurement studies report that
repeated runs of the same agentic task can differ by up to $30\times$ in total
tokens, and that agents predict their own consumption poorly
\citep{bai2026spend,lin2026bagen}. Unbounded consumption is listed among the
top risks for LLM applications \citep{owasp2025llm10}, and practitioners have
catalogued dozens of budget-overrun incidents in agent frameworks
\citep{khan2026tokenbudgets}.

A budget enforcer must decide, before each call, whether the call is
affordable. To do so it must plan for an unknown output length. Two plans are
common. Reserving the \emph{expected} output is permissive but gives no control
over how often the budget is exceeded. Reserving the \emph{maximum} output the
provider allows never exceeds the budget, but a capable model's maximum output
is often priced close to the whole budget, so the enforcer rarely admits it.
Neither plan lets the operator say how much risk they will accept.

We study the following guarantee. Given a run budget $B$, a tolerance
$\varepsilon$ and a confidence level $1-\delta$, deploy a controller $\pi$ such
that, with probability at least $1-\delta$ over calibration,
\begin{equation}\label{eq:goal}
  R(\pi) \;=\; \Pr_{\text{task},\ \text{sampling}}\bigl(\text{run cost under } \pi > B\bigr) \;\le\; \varepsilon .
\end{equation}
The controller is adaptive: it chooses models and admits or refuses calls based
on the money already spent, so per-step observations are not exchangeable and
standard per-call conformal guarantees do not compose directly into
\eqref{eq:goal}.

\paragraph{Contributions.}
\begin{enumerate}
  \item We formalize run-level budget risk for adaptive agent controllers in two
  operating modes: one where responses are never cut off and the risk is
  overrun, and one where a \texttt{max\_tokens} backstop prevents overrun and the
  risk is truncation (Section~\ref{sec:setup}).
  \item We show that per-step quantile bounds combined with a union bound or
  risk spending are valid but degenerate to worst-case reservation at practical
  calibration sizes, and explain why (Section~\ref{sec:perstep}).
  \item We calibrate a one-parameter policy family at the run level with
  Learn-then-Test \citep{angelopoulos2021ltt}, giving the guarantee
  \eqref{eq:goal} from i.i.d.\ runs with no distributional assumption on output
  lengths (Section~\ref{sec:method}). The statistical tool is not new; the
  contribution is the formulation and the evidence that run-level calibration
  is what makes the guarantee usable.
  \item We evaluate in a simulator with heavy-tailed outputs, growing context
  and adaptive routing, including a misspecified variant, and validate the
  guarantee over 500 independent calibrations (Section~\ref{sec:sim}). We
  then replay 2{,}100 real agent runs exactly under each policy
  (Section~\ref{sec:real}).
  \item We release the method as a reservation mode in the open-source
  Baar-Core library.
\end{enumerate}

% ---------------------------------------------------------------------------
\section{Related work}

\paragraph{Cost-aware routing and cascades.}
FrugalGPT \citep{chen2023frugalgpt} and RouteLLM \citep{ong2024routellm} trade
cost against quality by choosing among models, evaluated by average cost.
C3PO \citep{valkanas2025c3po} optimizes cascades under a probabilistic
constraint on the cost of a \emph{single query}. Our constraint is on the total
cost of a multi-step \emph{run} under an adaptive controller.

\paragraph{Budget enforcement for agents.}
\citet{khan2026tokenbudgets} catalogues 63 production overrun incidents and
enforces a dollar cap with an affine-typed Rust library, reporting static
over-reservation of 4--6$\times$. Green SARC \citep{besanson2026greensarc}
places enforcement gates in the agent loop and calibrates its per-action cost
forecasts with split-conformal prediction. These systems make the cap hard by
reserving conservatively, or calibrate individual actions; to our knowledge
none bounds the probability that an entire adaptive run exceeds its budget
while letting the operator trade that probability for capability.

\paragraph{Distribution-free risk control.}
Conformal prediction and its risk-control extensions
\citep{bates2021rcps,angelopoulos2022crc} bound the expected loss of a
calibrated predictor. Learn-then-Test \citep{angelopoulos2021ltt} recasts the
choice of a hyperparameter as multiple hypothesis testing and controls
arbitrary, non-monotone risks. Anytime-valid variants handle sequentially
arriving data \citep{hultberg2026anytime}, and ToolChain-CRC applies
trajectory-level conformal risk control to agent correctness under drift
\citep{opoku2026toolchaincrc}. DAPRO \citep{feldman2026spend} allocates a
compute budget dynamically across multi-turn LLM evaluations with
distribution-free coverage guarantees, satisfying the budget in expectation
rather than with high probability. We apply Learn-then-Test to budget risk, where
the relevant loss is defined over whole runs.

% ---------------------------------------------------------------------------
\section{Problem setup}\label{sec:setup}

A task $X$ is drawn i.i.d.\ from a distribution $\mathcal{D}$. A run proceeds
in steps $t = 1, 2, \dots$. Before step $t$ the controller knows the money spent
so far $s_{t-1}$, the prompt length $n_t$, and the history. It picks a model
$m_t$ and a step type $k_t$ (tool call or final answer), or refuses, which ends
the run without an answer. After the call it observes the output length $y_t$
and pays $c_t = p^{\text{in}}_{m_t} n_t + p^{\text{out}}_{m_t} y_t$.

\paragraph{Reservation policies.}
The controller admits a call if $s_{t-1} + p^{\text{in}}_m n_t +
p^{\text{out}}_m\, u(m,k) \le B$, where $u(m,k)$ is a planned output length. To
leave room to finish, a tool call is admitted only if a subsequent final answer
on the cheap model also fits. The plan $u$ defines the policy:
\begin{itemize}
  \item \textbf{estimate}: $u$ is the mean output length;
  \item \textbf{preflight}: $u = Y^{\max}_m$, the provider's maximum output;
  \item \textbf{calibrated} (ours): $u = \hat q_\beta(m,k)$, the empirical
  $(1-\beta)$ quantile of output length for stratum $(m,k)$, estimated on
  held-out logs. $\beta$ is a single scalar chosen by calibration.
\end{itemize}

\paragraph{Two modes.}
In \textbf{Mode A} responses are never cut off, and the run loss is
$L_A = \mathbf{1}\{\sum_t c_t > B\}$. In \textbf{Mode B} each call carries
\texttt{max\_tokens} equal to the largest output that still fits the remaining
budget. Overrun becomes impossible and the loss is truncation,
$L_B = \mathbf{1}\{\text{some response was cut off}\}$. Mode A is the default
in Baar-Core because silent truncation can corrupt downstream results without
any visible error. The risk of policy $\pi$ is $R(\pi) = \mathbb{E}[L(\pi)]$.

Preflight has $R = 0$ in both modes: it never overruns, and in Mode B its
\texttt{max\_tokens} is never binding.

% ---------------------------------------------------------------------------
\section{Why per-step bounds collapse}\label{sec:perstep}

The natural route to \eqref{eq:goal} bounds each step. Choose levels
$\alpha_t$ with $\sum_t \alpha_t \le \varepsilon$ (a union bound, or an online
``risk-spending'' schedule that spends a fraction of the remaining
$\varepsilon$ only when a step is tight), and plan $u_t$ at an upper
confidence bound on the $(1-\alpha_t)$ quantile. If every $y_t \le u_t$, the run
cannot overrun, so the overrun probability is at most $\sum_t \alpha_t$ plus
the calibration failure probability.

This is valid, but expensive in two ways.

\begin{proposition}\label{prop:n}
Let $Y_1, \dots, Y_n$ be i.i.d.\ and $Y_{(j)}$ their order statistics. The
smallest order statistic that is a valid $(1-\alpha)$ upper quantile bound with
confidence $1-\delta'$ is $Y_{(j)}$ with the least $j$ such that
$\Pr(\mathrm{Bin}(n, 1-\alpha) \ge j) \le \delta'$. No order statistic
qualifies, and the bound must fall back to $Y^{\max}$, unless
\[
  n \;\ge\; \frac{\log(1/\delta')}{-\log(1-\alpha)} \;\approx\; \frac{\log(1/\delta')}{\alpha}.
\]
\end{proposition}
\begin{proof}
$Y_{(j)}$ lies below the $(1-\alpha)$ quantile only if at least $j$ samples do,
which has probability at most $\Pr(\mathrm{Bin}(n,1-\alpha) \ge j)$. Even
$j = n$ requires $(1-\alpha)^n \le \delta'$.
\end{proof}

With $\varepsilon = 0.05$ spread over ten steps ($\alpha = 0.005$), and
$\delta'$ divided further across strata and quantile levels so that all bounds
hold simultaneously (e.g.\ $\delta' = 0.05/160$ for 4 strata and 40 levels),
the proposition requires about $1{,}600$ logged outputs per stratum merely to beat $Y^{\max}$. Risk
spending makes later steps stricter still, because the remaining budget of risk
shrinks geometrically.

The deeper problem is that the union bound charges every step for its own tail.
A single long response almost never causes a run-level overrun, because most
runs finish well under budget. The per-step bound for the capable model, at a
tail level of a fraction of a percent, is priced near the entire budget, so the
capable model is never admitted. Table~\ref{tab:valid} shows exactly this: the
per-step policy achieves zero overruns and zero use of the capable model,
indistinguishable from preflight. Run-level calibration avoids the union bound
altogether.

% ---------------------------------------------------------------------------
\section{Run-level calibration with Learn-then-Test}\label{sec:method}

Fix in advance a grid $\beta_1 < \beta_2 < \dots < \beta_J$ ordered from
conservative to aggressive, and the per-stratum quantile tables
$\hat q_\beta$, estimated on data disjoint from calibration. For each $\beta_j$
the calibrated policy $\pi_{\beta_j}$ is fully specified, so its run-level risk
$R_j = R(\pi_{\beta_j})$ is a fixed number. We test the null hypotheses
$H_j : R_j > \varepsilon$ in order.

\paragraph{Procedure.}
For $j = 1, 2, \dots$: execute $n$ fresh runs on i.i.d.\ tasks under
$\pi_{\beta_j}$ with budget $B$, count failures $F_j$, and compute the exact
binomial $p$-value
\[
  p_j \;=\; \Pr\bigl(\mathrm{Bin}(n,\varepsilon) \le F_j\bigr).
\]
If $p_j \le \delta$, reject $H_j$ and continue; otherwise stop. Deploy
$\hat\beta = \beta_{\hat\jmath}$, the last rejected level, or preflight if
$H_1$ is not rejected.

\begin{theorem}\label{thm:main}
If calibration runs are i.i.d.\ draws of (task, sampling randomness) from the
deployment distribution, then
$\Pr\bigl(R(\pi_{\hat\beta}) > \varepsilon\bigr) \le \delta$, where the
probability is over the calibration runs. This holds for any output-length
distributions, any dependence between steps, and with no monotonicity of $R_j$
in $j$.
\end{theorem}
\begin{proof}
If $R_j \le \varepsilon$ for all $j$, the deployed policy (any $\pi_{\beta_j}$,
or preflight with risk $0$) has risk at most $\varepsilon$. Otherwise let
$j^\ast$ be the smallest index with $R_{j^\ast} > \varepsilon$; it is
deterministic. The deployed policy has risk above $\varepsilon$ only if
$\hat\jmath \ge j^\ast$, which requires $H_{j^\ast}$ to be rejected. Under
$H_{j^\ast}$, $F_{j^\ast} \sim \mathrm{Bin}(n, R_{j^\ast})$ with
$R_{j^\ast} > \varepsilon$, which is stochastically larger than
$\mathrm{Bin}(n,\varepsilon)$; hence $\Pr(p_{j^\ast} \le \delta) \le \delta$.
\end{proof}

This is the fixed-sequence instance of Learn-then-Test
\citep{angelopoulos2021ltt}. Three properties matter for budgets. First, the
guarantee is about whole runs, so it accounts for adaptive routing, growing
context and correlated step lengths without modelling them. Second, the
per-stratum quantile tables only shape the policy family; if they are wrong,
the procedure selects a more conservative $\beta$ rather than losing validity.
Third, the same procedure controls truncation in Mode B by changing the loss.

\paragraph{Cost of calibration.}
Each tested level requires $n$ executed runs, because routing depends on
$\beta$ and runs cannot be replayed counterfactually. The total is
$(\hat\jmath + 1) n$ runs. In Mode A some calibration runs overrun by design;
this spend should be accounted for when budgeting the calibration itself.

% ---------------------------------------------------------------------------
\section{Simulation study}\label{sec:sim}

\paragraph{Simulator.}
Two models with list prices of $\$0.15/\$0.60$ (small) and $\$2.50/\$10.00$
(capable) per million input/output tokens, and maximum outputs of 4{,}096 and
8{,}192 tokens. Output lengths are lognormal per stratum (model $\times$ step
type), with the capable model and final answers longer. Each run starts with a
600-token prompt; each tool step appends the output and an observation of
50--800 tokens. The number of steps the task needs is geometric with mean 6,
capped at 30. The controller prefers the capable model while more than half
the budget remains, downgrades to the small model when the capable one does not
fit, and wraps up early on the small model if necessary. The run budget is
$B = \$0.02$. In the \textbf{misspecified} scenario every task has a hidden
difficulty that scales all its output lengths by $e^{d}$, $d \sim
\mathcal{N}(0, 0.6^2)$, so steps within a run are correlated and the
per-stratum model is wrong.

\paragraph{Settings.}
Quantile tables use 2{,}000 held-out samples per stratum. Learn-then-Test uses
$\beta \in \{0.001, 0.002, 0.005, 0.01, 0.02, 0.05, 0.1, 0.2, 0.3, 0.5\}$,
$n = 1{,}000$ calibration runs per level and $\delta = 0.05$. The per-step
baseline uses exact binomial quantile bounds with Bonferroni correction across
strata and quantile levels and lazy risk spending. Each policy is evaluated on
20{,}000 fresh test runs. ``Big'' is the fraction of executed steps on the
capable model, our proxy for how much capability the budget buys; ``done'' is
the fraction of runs that reach a final answer rather than being refused.

\begin{table}[t]
\centering
\caption{Well-specified simulator, $B=\$0.02$, $\delta=0.05$, 20{,}000 test
runs. Failure is overrun in Mode A and truncation in Mode B. Estimate and
preflight do not depend on $\varepsilon$.}
\label{tab:valid}
\begin{tabular}{llrrrr}
\toprule
Mode & Policy & $\hat\beta$ & Failure & Big & Done \\
\midrule
A & estimate & -- & 4.69\% & 44.6\% & 88.8\% \\
A & preflight & -- & 0 & 0 & 100\% \\
A & per-step, $\varepsilon \in \{.01,.05,.10\}$ & -- & 0 & 0 & 100\% \\
A & LTT, $\varepsilon = 0.01$ & 0.1 & 0.37\% & 31.3\% & 94.0\% \\
A & LTT, $\varepsilon = 0.05$ & 0.2 & 2.63\% & 41.4\% & 90.4\% \\
A & LTT, $\varepsilon = 0.10$ & 0.5 & 6.46\% & 45.2\% & 88.4\% \\
\midrule
B & estimate & -- & 13.55\% & 44.6\% & 88.8\% \\
B & LTT, $\varepsilon = 0.01$ & 0.01 & 0 & 0 & 100\% \\
B & LTT, $\varepsilon = 0.05$ & 0.05 & 3.08\% & 24.4\% & 96.3\% \\
B & LTT, $\varepsilon = 0.10$ & 0.1 & 5.26\% & 31.3\% & 94.0\% \\
\bottomrule
\end{tabular}
\end{table}

\paragraph{Results.}
Table~\ref{tab:valid} shows the well-specified case. Estimate mode overruns
4.7\% of runs regardless of the operator's tolerance, and in Mode B it
truncates 13.6\% of runs. Preflight and the per-step policy never fail but
never use the capable model. Learn-then-Test meets every target and, at
$\varepsilon = 0.05$, spends 41\% of steps on the capable model while
overrunning 2.6\% of runs, close to estimate mode's capability at about half
its overrun rate and with a certificate. At $\varepsilon = 0.01$ in Mode B no
aggressive level could be certified and the procedure correctly fell back to
near-preflight behaviour. The misspecified scenario gives the same picture
(Table~\ref{tab:broken}): estimate mode overruns 5.3\%, and Learn-then-Test
meets every target with 29--46\% capable-model steps.

\begin{table}[t]
\centering
\caption{Misspecified simulator (hidden per-task difficulty), Mode A.}
\label{tab:broken}
\begin{tabular}{lrrrr}
\toprule
Policy & $\hat\beta$ & Overrun & Big & Done \\
\midrule
estimate & -- & 5.25\% & 45.2\% & 85.9\% \\
preflight / per-step & -- & 0 & 0 & 100\% \\
LTT, $\varepsilon = 0.01$ & 0.1 & 0.50\% & 29.1\% & 92.4\% \\
LTT, $\varepsilon = 0.05$ & 0.2 & 2.65\% & 41.0\% & 87.6\% \\
LTT, $\varepsilon = 0.10$ & 0.5 & 7.05\% & 45.9\% & 85.7\% \\
\bottomrule
\end{tabular}
\end{table}

\paragraph{Is the guarantee real?}
Theorem~\ref{thm:main} is a statement about the calibration draw, so we repeat
the whole procedure 500 times per scenario with independent calibration runs
($\varepsilon = \delta = 0.05$, Mode A) and measure each deployed policy's
overrun rate on 4{,}000 fresh runs. The deployed policy violated
$\varepsilon$ in 1 of 500 calibrations in the well-specified scenario (0.2\%;
95\% Clopper--Pearson interval 0.0--1.1\%) and 5 of 500 in the misspecified
scenario (1.0\%; 0.3--2.3\%), both well below $\delta = 5\%$. These counts are
slightly inflated by Monte Carlo error in the 4{,}000-run estimate; the largest
measured overrun rate was 5.17\%. Mean overrun was 3.0\% and 2.8\%, and
capable-model share was 42.4\% and 40.4\%.

\paragraph{The price of control.}
Learn-then-Test completes fewer runs than preflight (88--94\% versus 100\%),
because a more aggressive plan admits capable-model steps that later leave too
little budget to finish, and the controller refuses. Completion is not part of
the controlled loss here. A loss of ``overrun \emph{or} incomplete'' is
equally valid under Theorem~\ref{thm:main} and would trade capability for
completion; we leave its study to future work.

% ---------------------------------------------------------------------------
\section{Real-model study}\label{sec:real}

\paragraph{Setup.} We ran a ReAct-style agent \citep{yao2023react} on the
HotpotQA distractor validation set \citep{yang2018hotpotqa}. The agent sees the
titles of the ten context paragraphs, may call a \texttt{lookup(title)} tool
to read one, and stops with a final answer or after eight steps (temperature
0.7, answers scored by token F1). We collected 1{,}500 runs with gpt-4o-mini
and 600 with gpt-4o on the same shuffled tasks, logging every call's prompt
tokens, Baar-Core's pre-call prompt estimate, output tokens and cost. Total
spend was \$1.85, itself hard-capped by Baar-Core in preflight mode.

\paragraph{Exact replay.} A controller that only admits or refuses calls does
not change what the agent does before it refuses. Each logged run can therefore
be replayed under any reservation policy and budget, giving the exact outcome
that policy would have had on that task and sampling draw. Reservations use
the pre-call prompt estimate that was logged before each call, exactly as the
controller saw it. We set the run
budget $B$ to the 50th, 75th or 90th percentile of unconstrained run cost, so
that it binds for 50\%, 25\% or 10\% of runs. For Learn-then-Test, a random
quarter of each model's runs fits the quantile tables; the rest form a
population whose true risk under every $\beta$ is computed exactly. We draw
calibration sets of half the population size i.i.d.\ from it, run the
procedure ($\varepsilon = \delta = 0.05$), and repeat 500 times.

\paragraph{What real traces look like.} Outputs are short and light-tailed:
the 99th percentile is 50--67 tokens for both models and both step types, and
the maximum observed is 181. The provider maximum of 16{,}384 tokens is
250--330 times the 99th percentile. Cost is dominated by input tokens, as
reported for coding agents \citep{bai2026spend}; still, the most expensive run
cost 8--11 times the median. The logged pre-call prompt estimate
\emph{under}-counts in 94\% of calls, by a median of 29 tokens (9\%; 37
tokens on the first step), because Baar-Core's token counter could not take
the tool schema, which providers bill as prompt tokens. Recounting with the
schema matched the billed prompt to within one token. This study exposed the
bug, and it is fixed in Baar-Core; we keep the original estimate in the replay
because it is what the controller actually used. This matters below.

\begin{table}[t]
\centering
\caption{Real-trace replay on HotpotQA ($\varepsilon=\delta=0.05$). $B$ is set
at a percentile of unconstrained run cost. ``Done'' is the fraction of runs
that finish; F1 counts refused runs as 0. LTT shows the most frequently
selected $\hat\beta$ and its exact risk; ``viol.'' is the number of the 500
calibrations whose selected policy had risk above $\varepsilon$. Preflight
finished no run at any of these budgets.}
\label{tab:real}
\small
\begin{tabular}{llrrrrrrrr}
\toprule
& & \multicolumn{3}{c}{estimate} & \multicolumn{5}{c}{Learn-then-Test} \\
\cmidrule(lr){3-5}\cmidrule(lr){6-10}
Model & $B$ & Overrun & Done & F1 & $\hat\beta$ & Overrun & Done & F1 & viol. \\
\midrule
4o-mini & p50 & 5.3\% & 54.7\% & .328 & 0.1 & 2.8\% & 51.0\% & .308 & 1/500 \\
4o-mini & p75 & 1.9\% & 77.0\% & .489 & 0.5 & 2.8\% & 78.0\% & .496 & 0/500 \\
4o-mini & p90 & 0.4\% & 89.7\% & .584 & 0.5 & 0.6\% & 90.0\% & .586 & 0/500 \\
4o & p50 & 4.9\% & 54.7\% & .317 & 0.05 & 1.6\% & 44.4\% & .258 & 1/500 \\
4o & p75 & 1.8\% & 75.1\% & .458 & 0.3 & 1.6\% & 74.4\% & .454 & 0/500 \\
4o & p90 & 0.4\% & 88.4\% & .554 & 0.5 & 1.1\% & 89.3\% & .559 & 0/500 \\
\bottomrule
\end{tabular}
\end{table}

\paragraph{Results.} Table~\ref{tab:real} reports the replay. Four findings
stand out.
\begin{enumerate}
  \item \textbf{Worst-case reservation is unusable.} At every budget, the
  reserved cost of a 16{,}384-token output exceeds the budget, so preflight
  finishes no run. This is the failure the method exists to fix.
  \item \textbf{The guarantee holds on real runs.} Across six settings and
  3{,}000 calibrations, at most one calibration per setting selected a policy
  with risk above $\varepsilon$, against an allowed $\delta = 5\%$.
  \item \textbf{Outcomes match estimate mode; the gain is the certificate.} At
  the 75th and 90th percentile budgets, estimate mode happened to overrun
  less than $\varepsilon$, and the certified policy reaches the same
  completion and F1. At the tight median budget estimate mode sits at
  4.9--5.3\%, at the edge of the tolerance; Learn-then-Test cuts overruns to
  1.6--2.8\% at a cost of 4--10 points of completion. The difference is that
  the certified policy's risk is guaranteed in advance, while estimate mode's
  depends on a workload the operator cannot see beforehand.
  \item \textbf{Residual risk has more than one source.} Even at
  $\beta = 0.001$, gpt-4o overruns 0.4--0.7\% of runs at the tighter budgets.
  Replaying with the true prompt counts separates two causes. At the median
  budget all of these overruns remain, so they come from outputs longer than
  any in the 150-run quantile table (whose extreme quantile is 67 tokens,
  against an observed maximum of 157). At the 75th percentile budget,
  correcting the prompt count removes two thirds of them (0.7\% to 0.2\%).
  Run-level calibration absorbs both errors because it measures realized
  overruns; a per-step output bound covers neither.
\end{enumerate}

Two caveats. On this task the capable model was not more accurate (F1 0.61
versus 0.63 unconstrained), so we evaluate each model separately rather than
the cheap/capable routing of Section~\ref{sec:sim}; routing changes the
trajectory and cannot be replayed exactly. And HotpotQA outputs are short, so
the gap between preflight and calibrated reservation is mostly about the
provider maximum; agents that write code or long answers would have heavier
tails.

% ---------------------------------------------------------------------------
\section{Implementation}

The calibrated plan is a reservation mode in Baar-Core, an open-source Python
library that enforces LLM budgets before each request is sent. Baar-Core
estimates the cost of a call locally, atomically reserves it against a
budget, reconciles with the provider-reported cost afterwards, and routes
between a cheap and a capable model with downgrading. It already offers the
estimate and preflight plans as \texttt{estimate} and \texttt{preflight}
output modes, and Mode B as an opt-in \texttt{truncate} mode.

The method adds a \texttt{calibrated} mode and a dependency-free module
\texttt{baar.calibration}. \texttt{QuantileTable} builds per-stratum output
quantiles from logged calls; \texttt{ltt\_select} runs the fixed-sequence
procedure of Section~\ref{sec:method} given a function that executes $n$ runs
at a level $\beta$ and returns the failure count; exact binomial $p$-values are
computed in log space. The router then reserves the calibrated quantile:
\begin{verbatim}
table = QuantileTable.from_records(logged_calls)
res = ltt_select(betas, run_agent_batch, n=1000, eps=0.05, delta=0.05)
router = BAARRouter(budget=0.10, on_output_overrun="calibrated",
                    output_quantiles=table, calibration_beta=res)
\end{verbatim}
The mode never sets \texttt{max\_tokens}, and it falls back to preflight when
no level is certified or a stratum has too few samples. One gap remains: the
router does not yet know whether a step is a tool call or a final answer, so
it pools both step types per model; per-type strata are available through the
lower-level \texttt{plan\_output} call. The code is available at commit
\texttt{c989f34b4498} of the repository below and will ship in the next
Baar-Core release.

% ---------------------------------------------------------------------------
\section{Limitations}

\begin{itemize}
  \item \textbf{Exchangeability.} The guarantee assumes deployment tasks come
  from the calibration distribution. Price changes, model updates or a new
  task mix require recalibration; anytime-valid monitoring
  \citep{hultberg2026anytime} is a natural complement.
  \item \textbf{Calibration cost.} Each tested level costs $n$ real runs at
  budget $B$. With $n = 1{,}000$ and up to ten levels this is affordable for
  cheap agents and expensive for long-horizon ones. Smaller $n$ forces more
  conservative choices.
  \item \textbf{Coarse policy family.} One scalar $\beta$ shared by all strata
  is simple to certify but not optimal; richer families need multiple-testing
  corrections that cost power.
  \item \textbf{Completion.} As discussed, completion falls as capability
  rises unless it is included in the loss.
  \item \textbf{Scope of the real study.} One task family, two models from
  one provider, and single-model runs; the routing results are simulated. On
  real traces the certified policy matched, rather than beat, expected-cost
  reservation in outcomes. Exact replay covers admit/refuse controllers only.
\end{itemize}

% ---------------------------------------------------------------------------
\section{Conclusion}

Operators of LLM agents can now state a budget tolerance instead of choosing
between unbounded risk and worst-case reservation. Per-step tail bounds are
the obvious way to get such a guarantee, but they spend the risk tolerance on
tail events that rarely matter. Calibrating whole runs with Learn-then-Test
gives a finite-sample guarantee from i.i.d.\ runs alone and, in simulation,
recovers most of the capability that worst-case reservation gives up.

\paragraph{Reproducibility.} Simulator, calibration code, the real-model
harness and all raw results are in the \texttt{research/} directory of
\url{https://github.com/orvi2014/Baar-Core}, at commit
\href{https://github.com/orvi2014/Baar-Core/tree/c989f34b44983407b5a4d4f9f1d2108eb4d54e7d}{\texttt{c989f34b4498}}.

\paragraph{AI usage.} Generative AI tools (Claude Opus and Claude Sonnet,
Anthropic, through Claude Code) were used to write simulation and analysis
code and to draft and edit this paper. The author chose the research question
and method, checked the proofs, reviewed all code and text, and is
responsible for the content. All reported numbers are produced by the released
scripts.

\bibliographystyle{plainnat}
\bibliography{refs}
\end{document}
